\section{Keep two histories while retaining interleaved output}
For each block, gather left codes in chronological order into one window and right
codes into another. Each window has its own seven-value history. Both use the same
eight coefficients, but no sample from one channel enters the other channel's sum.
After filtering, place the left and right results back into their original frame
positions. Update each history from its own input window.

\begin{figure}[htbp]
\centering
\begin{tikzpicture}[x=2.1cm,y=0.7cm,alt={Eight consecutive signed sixteen-bit sample slots form four stereo frames. Slots from left to right are L0,R0,L1,R1,L2,R2,L3,R3. Left filtering follows slots zero,two,four,six; right filtering follows slots one,three,five,seven. Each slot is two bytes and each complete frame is four bytes.}]
\foreach \frame in {0,1,2,3}{
 \pgfmathtruncatemacro{\leftslot}{2*\frame}
 \pgfmathtruncatemacro{\rightslot}{2*\frame+1}
 \draw[fill=black!8] (\leftslot*0.5,0) rectangle +(0.5,1);
 \draw (\rightslot*0.5,0) rectangle +(0.5,1);
 \node at (\leftslot*0.5+0.25,0.5) {$L_{\frame}$};
 \node at (\rightslot*0.5+0.25,0.5) {$R_{\frame}$};
 \node[below] at (\frame+0.5,-0.15) {frame \frame};
}
\end{tikzpicture}
\caption{Interleaved storage contains two separate time sequences. The slot order is also specified in the figure description; this is a native drawing, not an IDE capture.}
\end{figure}

The stride from left frame zero to left frame one is four bytes, not two. A pointer
increment that is correct for mono therefore selects the wrong sequence in stereo.
Use a left-only impulse and then a right-only impulse to expose channel mixing. Also
choose different signals in the two channels: an error that swaps identical channels
cannot be detected by comparing those identical outputs.
